\section{Step 10 Measure this boundary before adding search}\label{h263-sad-s10}
Count completed instructions from function entry through its return to a stopped
caller continuation. Include admission checks and reduction; exclude preparation
of the caller's data. Record which implementation and input produced the count.
An instruction count is reproducible simulator work, while host wall time also
includes scheduling and runtime effects.

The checked reference used 3245 completed instructions in scalar mode and 207 in
MSA mode for each admitted worked block. A second valid MSA implementation used
239, because it formed magnitudes with separate maximum, minimum and subtraction
operations. All produced the same exact sums and active-row accesses. Neither
207 nor the ratio between these counts is a grading threshold.

\paragraph{Observe} Record counts for equal, maximum-difference and worked
windows. Confirm correctness independently before reporting the counts.
\Needspace{8\baselineskip}
\paragraph{Check} What can these measurements establish by themselves?
\begin{itemize}
\item[A.] Completed instruction work for the named kernel boundaries and inputs.
\item[B.] The running time of the finished video encoder.
\item[C.] A requirement that every correct solution execute exactly 207 instructions.
\end{itemize}
